\documentclass{gamadays}
\usepackage{hyperref}

% ------------------------------------------
% TITLE
% ------------------------------------------

\title{\textbf{SImulating online deBAtes (SIMBA)}}


% ------------------------------------------
% AUTHOR(S)
% ------------------------------------------

% 1 auteur
% \author{L. Author \\
%  Employer, laboratory acronym}
% \date{email}
 
% plusieurs auteurs 
\author{A. Gregor Stolte\up{1}, B. Nicolas Darcel\up{1}, C. Patrick Taillandier\up{2}\\[6pt]
\up{1} AgroParisTech, CRNH-IdF, UMR914 Nutrition Physiology and Ingestive Behavior, Paris, France\\
\up{2} MIAT, University of Toulouse, INRA, 31320, Castanet-Tolosan, France\\
...}

\date{Nicolas.darcel@agroparistech.fr}

\begin{document}

\maketitle


\begin{keywords}
Agent-Based Simulation, Food Decision-Making, Social Influence
\end{keywords}

\begin{abstract}
Understanding how we eat and with whom is critical to a sustainable future in terms of food and biodiversity. The choice of what to eat has profound impacts on our health, the environment and the Greenhouse Gas Emissions (GHGs) resulting from our food choices. Existing research in this domain has primarily focused on behavioral interventions grounded in social norms and nudges, yielding mixed results. A proposed shift from this is investigating deliberation — a structured process facilitating the exchange of diverse attitudes — as a lens through which to understand how individuals' positions on meat consumption evolve through group interaction.\newline

We present SIMBA, an open-source Agent-Based Model developed in GAMA that simulates the functional dynamics of deliberation between multiple individuals. The model was designed to reproduce pre- and post-deliberation attitude shifts observed across 55 group debates (N = 459) from an empirical study on collective eating behavior and red meat consumption reduction (Dheilly et al., unpublished). The goal was to better understand the mechanisms by which individuals influence each other in collective eating settings (e.g., restaurants, university cafeterias). SIMBA draws on three established models of social influence — consensus (agents gradually converge), clustering (agents only influence those with similar views) and bipolarization (dissimilar agents actively push apart) — to recreate individual attitude trajectories observed in the empirical debate context. Agents are initialized solely with pre-deliberation attitudes, ensuring that any predictive improvement stems from the interaction dynamics rather than from additional individual-level covariates. \newline

Model calibration was performed on an 80/20 stratified train-test split of the debates, with an explicit held-out set reserved for validation. The two-step calibration process used Latin-Hypercube Sampling (200 samples) to explore the parameter space, followed by a Genetic Algorithm optimizing Mean Absolute Error (MAE) relative to empirical attitude shifts to identify final parameter values applied to the held-out debates. \newline

Results indicate that none of the three social influence variants produced significantly lower MAE than a naive no-change baseline at the individual, debate or global level. The Genetic Algorithm converged toward near-zero predicted attitude change rather than learning genuine debate dynamics, suggesting that simple opinion-averaging and bounded-confidence interaction rules are insufficient to capture the deliberative processes observed empirically. These findings establish the baseline for the second phase of this study, which integrates argumentation dynamics — specifically Dung-style argument attack and defense relations — to test whether richer interaction mechanisms improve predictive performance.
\end{abstract}

\begin{additionnalMaterial}
\href{https://osf.io/nfue6/overview}{OSF Protocol}
\end{additionnalMaterial}

\end{document}

